\section*{Problems}

\subsection*{Problem Statements}

% Remember
\begin{problema}
\label{prob:c7c2f12}
A financial analyst fits a time-series regression with $n=10$ consecutive observations and obtains temporally ordered residuals
\begin{align}
e=\{2.1,\ 1.8,\ -0.5,\ -1.2,\ -2.0,\ -1.5,\ 0.8,\ 1.9,\ 2.2,\ 1.0\}.
\end{align}
Calculate Durbin--Watson $DW=\dfrac{\sum_{i=2}^n(e_i-e_{i-1})^2}{\sum_{i=1}^n e_i^2}$ and determine whether residuals show positive autocorrelation, negative autocorrelation, or no autocorrelation.
\end{problema}

% Understand
\begin{problema}
\label{prob:90810fe}
In a regression of sales ($Y$) on advertising expenditure ($X$), residuals are plotted against fitted values. For low fitted values ($\hat Y<50$), residuals range over $\pm2$; for intermediate values ($50\le\hat Y<150$), over $\pm8$; and for high values ($\hat Y\ge150$), over $\pm20$.
\begin{enumerate}
\item Identify the geometric pattern and the classical assumption it violates.
\item Propose a corrective transformation and justify why it would be effective here.
\end{enumerate}
\end{problema}

% Apply
\begin{problema}
\label{prob:46ca545}
A credit analyst builds a model with $p=3$ potentially correlated financial predictors. Regressing each predictor on the other two gives auxiliary coefficients of determination $R_1^2=0.85$, $R_2^2=0.55$, and $R_3^2=0.30$.
\begin{enumerate}
\item Calculate $\text{VIF}_j=1/(1-R_j^2)$ for each predictor.
\item Using the conventional critical threshold $\text{VIF}>5$, determine which predictor(s) require intervention.
\item Propose and justify a remedy (Ridge, Lasso, or variable removal) for the most problematic predictor.
\end{enumerate}
\end{problema}

% Analyze
\begin{problema}
\label{prob:269eee0}
\textbf{Derivation of Cook's distance from its matrix definition:} starting from
\begin{align}
D_i=\frac{(\hat{\boldsymbol\beta}-\hat{\boldsymbol\beta}_{(-i)})^T(\mathbf X^T\mathbf X)(\hat{\boldsymbol\beta}-\hat{\boldsymbol\beta}_{(-i)})}{(p+1)\hat\sigma^2},
\end{align}
and using the Allen--PRESS theorem (which expresses $\hat{\boldsymbol\beta}-\hat{\boldsymbol\beta}_{(-i)}$ in terms of ordinary residual $e_i$, leverage $h_{ii}$, and row $\mathbf x_i$), derive the simplified computational formula
\begin{align}
D_i=\frac{r_i^2}{p+1}\cdot\frac{h_{ii}}{1-h_{ii}},
\end{align}
where $r_i$ is the internally studentized residual. Prove why observations with $h_{ii}\to1$ produce explosive Cook distances regardless of the ordinary residual magnitude.
\end{problema}

% Evaluate
\begin{problema}
\label{prob:dc2ca7b}
In a model with $n=50$ observations and $p=4$ predictors (5 parameters including the intercept), three suspicious observations have leverage ($h_{ii}$) and internally studentized residual ($r_i$):
\begin{align}
\text{Obs. A: }h_{ii}=0.15,\ r_i=1.2;\qquad \text{Obs. B: }h_{ii}=0.45,\ r_i=0.8;\qquad \text{Obs. C: }h_{ii}=0.08,\ r_i=3.5.
\end{align}
\begin{enumerate}
\item Calculate Cook's distance for each observation using $D_i=\dfrac{r_i^2}{p+1}\cdot\dfrac{h_{ii}}{1-h_{ii}}$.
\item Using the conventional threshold $4/n$, determine which observations are influential.
\item Evaluate the distinct mechanism by which B and C are influential (high leverage versus large residual), even though A is not despite moderate values of both quantities.
\end{enumerate}
\end{problema}

% Create
\begin{problema}
\label{prob:f0b580a}
Design an original heteroscedasticity-diagnostics problem (different from the advertising-sales case): an industrial-energy model uses plant size and $p=4$ other process variables, fitted to $n=80$ observations. An auxiliary Breusch--Pagan regression of squared residuals on the original predictors gives $R_{\text{aux}}^2=0.22$. Calculate $LM=nR_{\text{aux}}^2$, compare it with $\chi^2_{0.05,4}=9.488$, and decide whether there is evidence of heteroscedasticity at $\alpha=0.05$.
\end{problema}

\subsection*{Detailed Solutions}

% Remember
\begin{solproblema}[prob:c7c2f12]
The numerator is $\sum_{i=2}^{10}(e_i-e_{i-1})^2=14.79$. The denominator is $\sum_{i=1}^{10}e_i^2=25.68$. Therefore $DW=14.79/25.68\approx0.576$. Since $DW\ll2$, there is strong evidence of \textbf{positive autocorrelation}: consecutive residuals tend to have the same sign, violating independence.
\end{solproblema}

% Understand
\begin{solproblema}[prob:90810fe]
\begin{enumerate}
\item Residual amplitude grows systematically with the fitted value ($\pm2\to\pm8\to\pm20$), producing the classic \textbf{funnel shape} in a residual-versus-fitted plot. This violates \textbf{homoscedasticity}: error variance is not constant but increases with response level.
\item Heteroscedasticity with variance approximately proportional to the mean (or its square) is typically corrected by applying $\log(Y)$ to the response. The transformation compresses large absolute differences more than small ones and stabilizes relative residual spread.
\end{enumerate}
\end{solproblema}

% Apply
\begin{solproblema}[prob:46ca545]
\begin{enumerate}
\item $\text{VIF}_1=1/(1-0.85)\approx6.667$, $\text{VIF}_2=1/(1-0.55)\approx2.222$, and $\text{VIF}_3=1/(1-0.30)\approx1.429$.
\item Only predictor 1 exceeds the threshold ($\text{VIF}_1\approx6.667>5$), indicating severe multicollinearity; predictors 2 and 3 are tolerable.
\item Since predictor 1 may have theoretical relevance, \textbf{Ridge regression} is most appropriate: it continuously penalizes the $L_2$ norm and stabilizes $(\mathbf X^T\mathbf X+\lambda\mathbf I)^{-1}$ without discarding variables.
\end{enumerate}
\end{solproblema}

% Analyze
\begin{solproblema}[prob:269eee0]
By the Allen--PRESS theorem, deleting observation $i$ changes the coefficient vector by
\begin{align}
\hat{\boldsymbol\beta}-\hat{\boldsymbol\beta}_{(-i)}=\frac{(\mathbf X^T\mathbf X)^{-1}\mathbf x_i e_i}{1-h_{ii}}.
\end{align}
Substitution into Cook's quadratic form gives
\begin{align}
D_i=\frac{e_i^2}{(p+1)\hat\sigma^2}\cdot\frac{h_{ii}}{(1-h_{ii})^2},
\end{align}
using $\mathbf x_i^T(\mathbf X^T\mathbf X)^{-1}\mathbf x_i=h_{ii}$. With internally studentized residual $r_i=e_i/(\hat\sigma\sqrt{1-h_{ii}})$, so $e_i^2/\hat\sigma^2=r_i^2(1-h_{ii})$,
\begin{align}
D_i=\frac{r_i^2}{p+1}\cdot\frac{h_{ii}}{1-h_{ii}}.
\end{align}
The factor $h_{ii}/(1-h_{ii})$ diverges as $h_{ii}\to1$, regardless of $r_i$: an observation at an extreme predictor position can almost determine the fitted line, and the same vanishing $(1-h_{ii})$ term masks the discrepancy until the point is removed.
\end{solproblema}

% Evaluate
\begin{solproblema}[prob:dc2ca7b]
With $p+1=5$,
\begin{align}
D_A=\frac{1.44}{5}(0.1765)\approx0.0508,\qquad D_B=\frac{0.64}{5}(0.8182)\approx0.1047,\qquad D_C=\frac{12.25}{5}(0.0870)\approx0.2130.
\end{align}
The threshold is $4/50=0.08$: A is not influential, while B and C are influential. B is influential through \textbf{high leverage} ($h_{ii}=0.45$) despite a moderate residual; its unusual predictor position can pull the fitted line. C is influential through a \textbf{large residual} ($r_i=3.5$, a possible outlier) despite low leverage. A has neither mechanism at a severe level.
\end{solproblema}

% Create
\begin{solproblema}[prob:f0b580a]
\begin{align}
LM=nR_{\text{aux}}^2=80\times0.22=17.6.
\end{align}
Since $17.6>\chi^2_{0.05,4}=9.488$, \textbf{reject $H_0$}: there is significant evidence of \textbf{heteroscedasticity} in the industrial-energy model at the 5\% level. OLS estimators remain unbiased but are no longer efficient, and reported standard errors may be incorrect; use Huber--White robust errors or transform the response.
\end{solproblema}
